\documentclass[11pt]{article}
\usepackage[a4paper,margin=1in]{geometry}
\usepackage{amsmath,amssymb,amsthm}
\usepackage{mathtools}
\usepackage{xcolor}
\usepackage{enumitem}
\usepackage{microtype}
\usepackage[colorlinks,linkcolor=blue!50!black,citecolor=blue!50!black,urlcolor=blue!50!black]{hyperref}
\usepackage{cleveref}

% Theorem-like environments share one counter, numbered within sections.
\theoremstyle{plain}
\newtheorem{theorem}{Theorem}[section]
\newtheorem{lemma}[theorem]{Lemma}
\newtheorem{proposition}[theorem]{Proposition}
\newtheorem{corollary}[theorem]{Corollary}
\theoremstyle{definition}
\newtheorem{definition}[theorem]{Definition}
\newtheorem{example}[theorem]{Example}
\theoremstyle{remark}
\newtheorem{remark}[theorem]{Remark}

% Handy shortcuts. Add your own notation here.
\newcommand{\R}{\mathbb{R}}
\newcommand{\N}{\mathbb{N}}
\DeclarePairedDelimiter\abs{\lvert}{\rvert}
\DeclarePairedDelimiter\norm{\lVert}{\rVert}

\title{A Note on Contractive Maps in Weighted Sequence Spaces}
\author{First Author\thanks{Department of Mathematics, Example University. \texttt{first@example.edu}}
  \and Second Author\thanks{School of Sciences, Sample Institute.}}
\date{\today}

\begin{document}
\maketitle

\begin{abstract}
State the main theorem informally and explain why it is interesting. Mention the key technique in one sentence.
\end{abstract}

\noindent\textbf{MSC 2020:} 47H10, 46B45.\quad \textbf{Keywords:} fixed point, weighted $\ell^p$ space, contraction.

\section{Introduction}
Give background, state the main result early and outline the paper.
Our main result is \cref{thm:main}, which builds on the classical argument in \cite{banach}.

\section{Preliminaries}
\begin{definition}\label{def:weighted}
Let $w = (w_n)_{n\in\N}$ be a sequence of positive weights. The weighted space $\ell^p_w$ consists of real sequences $x$ with
\[
  \norm{x}_{p,w} \coloneqq \Bigl( \sum_{n=1}^{\infty} w_n \abs{x_n}^p \Bigr)^{1/p} < \infty .
\]
\end{definition}

\begin{lemma}\label{lem:complete}
For every $1 \le p < \infty$ the space $(\ell^p_w, \norm{\cdot}_{p,w})$ is complete.
\end{lemma}

\begin{proof}
Replace this sketch with your argument. Take a Cauchy sequence, show coordinatewise convergence, then use Fatou's lemma to bound the norm of the limit.
\end{proof}

\section{Main result}
\begin{theorem}\label{thm:main}
Let $T\colon \ell^p_w \to \ell^p_w$ satisfy $\norm{Tx - Ty}_{p,w} \le q \norm{x-y}_{p,w}$ for some $q \in [0,1)$. Then $T$ has a unique fixed point.
\end{theorem}

\begin{proof}
By \cref{lem:complete} the space is complete. Fix $x_0$ and set $x_{k+1} = T x_k$. Then
\begin{equation}\label{eq:geometric}
  \norm{x_{k+1} - x_k}_{p,w} \le q^k \norm{x_1 - x_0}_{p,w},
\end{equation}
so $(x_k)$ is Cauchy by \eqref{eq:geometric}. Its limit is a fixed point, and uniqueness follows because two fixed points $x,y$ give $\norm{x-y} \le q\norm{x-y}$.
\end{proof}

\begin{corollary}
Every iterate sequence of $T$ converges to the fixed point at a geometric rate.
\end{corollary}

\begin{example}
Take $w_n = 2^{-n}$ and $(Tx)_n = \tfrac{1}{2} x_{n+1}$. Work through the constants here.
\end{example}

\begin{remark}
Use remarks for comments that are not essential to the logical flow.
\end{remark}

\section{Open questions}
\begin{enumerate}[label=(Q\arabic*)]
  \item Does the result extend to $p = \infty$?
  \item What happens for non-summable weights?
\end{enumerate}

\section*{Acknowledgements}
Thank colleagues and grant bodies here.

\begin{thebibliography}{9}
\bibitem{banach} S.~Banach, Sur les op\'erations dans les ensembles abstraits et leur application aux \'equations int\'egrales, \emph{Fund. Math.} 3 (1922), 133--181.
\bibitem{book} A.~Author, \emph{Topics in Functional Analysis}, Example Publishers, 2015.
\end{thebibliography}

\end{document}
